% **************************************************************************************************************
% Thesis-Vorlage der Hochschule Furtwangen, Digitale Medien
%
% Basierend auf der Vorlage von Patrick Heneise; 2010 erweitert und angepasst von Martin Weber
% **************************************************************************************************************

\documentclass[a4paper,12pt,twoside,openright,numbers=noenddot]{scrreprt}
\addtokomafont{disposition}{\rmfamily}

% Packages
% ********************************************************************
\usepackage[utf8]{inputenc}
\usepackage[T1]{fontenc} % SUPER WICHTIG! Läd "neue" Schriften: Schönere Zeichen (»«), korrekte (!) Generierung von Sönderzeichen (wichtig bei copy&paste-Export: Der schlug sonst immer fehl). Ermöglicht außerdem korrekte Silbentrennung von Wörtern mit Ümläüten.
\usepackage{textcomp} % Interrobang ftw!
\usepackage[ngerman]{babel}
\usepackage[pdftex]{graphicx}
\usepackage{xcolor}
\usepackage{scrpage2} 
\usepackage[hyphens]{url}
\usepackage{natbib}
\usepackage{listings}
\usepackage{makeidx}
\usepackage{hyperref}
\usepackage[final]{pdfpages}
\usepackage{eurosym} % Fuer das Euro-Symbol mit \euro{}
\usepackage[inner=3.2cm,outer=2.2cm,top=2cm,bottom=2cm,includeheadfoot]{geometry}
\makeindex
\setcounter{tocdepth}{3} % Tiefe des Inhaltsverzeichnisses (3; bis zu 1.1.1 runter)
\setcounter{secnumdepth}{3} % Tiefe der Nummerierung innerhalb des Dokuments (3; bis zu 1.1.1 runter)
\setcitestyle{aysep={,},citesep={;}}
\usepackage{subfigure}  % Hiermit lassen sich mehrere Grafiken beinander anordnen,welche dann fortlaufend mit a, b usw. bezeichnet werden.
\usepackage{amsmath} % Für die mathematische Darstellung der Formeln zur Preisberechnung
\usepackage{listings} % Für die Quellcode-Darstellung, s. http://blog.robert-kummer.de/2006/04/latex-quellcode-listing.html
\usepackage{multirow} % Für zusammenfassung von Reihen in Tabellen; s. http://andrewjpage.com/index.php?/archives/43-Multirow-and-multicolumn-spanning-with-latex-tables.html
\usepackage{rotating} % Zur Rotation von Elementen
\clubpenalty = 10000 \widowpenalty = 10000 \displaywidowpenalty = 10000 %Gegen Schusterjungen und Hurenkinder


\usepackage{chngcntr} % Dient zur durchlaufenden Nummerierung der Abbildungen. Der Standard von scrreport ist [Chapter].[Nummer]
\counterwithout{figure}{chapter}
\counterwithout{table}{chapter}


\usepackage{remreset} %Dient dazu, die Fussnoten im kpl. Dokument durchnummerieren zu können. Ansonsten fangen sie bei jedem Chapter wieder bei 1 an.
\makeatletter
\@removefromreset{footnote}{chapter}
\makeatother


\hypersetup{ %Alle URLS schwarz färben
	colorlinks,%
	citecolor=black,%
	filecolor=black,%
	linkcolor=black,%
	urlcolor=black
}

\renewcommand{\harvardurl}[1]{\small{\url{#1}}}
\renewcommand{\dictumwidth}{0.6\textwidth} % Ändert die Breite des Zitats direkt nach der Titelseite.
\setkomafont{dictum}{\rmfamily \large} %Serifentext und Größefür das Zitat


\KOMAoptions{cleardoublepage=empty}
% Zeugs am Anfang: Titel, Abstract, Absichtserklärung, Inhaltsverzeichnis, Abbildungsverzeichnis.
% ********************************************************************
\begin{document}
\shorthandoff{"} % Verhindert, dass Anführungszeichen nach einem Vokal schief laufen.
\pagestyle{empty}

 \thispagestyle{empty}
\begin{titlepage}

\begin{figure}
\begin{flushright}
 \includegraphics[width=0.3\textwidth]{./Images/200px-Seal_Leiden_University.jpg}
 \end{flushright}
\end{figure}

\begin{flushleft} 
% Title
{\small \textsf{\textbf{Masterthesis}}} \\[0.5cm]
{\Huge \textsf{\textbf{Your Awesome Thesis Title}}} \\[0.0cm]
{\Large \textsc{And Some Cool Subtitle}} \\[1cm]
 
% Author and supervisor
John Doe\\
john@doe.com\\[1cm]

\color{gray}Supervisor: Prof. Dr. Superprof\\
\color{gray}Universiteit Leiden\\[1cm]

\vfill 
\end{flushleft}
 
\end{titlepage}


\cleardoublepage

\thispagestyle{empty}
\vspace*{5cm}
\dictum[Schlauer Autor, 1900--2010]{Hier steht ein schlauer Spruch.} %Zitat zu Beginn

\pagestyle{plain}
 \begin{abstract}
Lorem ipsum dolor sit amet, consectetuer adipiscing elit, sed diam nonummy nibh euismod tincidunt ut laoreet dolore magna aliquam erat volutpat. Ut wisi enim ad minim veniam, quis nostrud exerci tation ullamcorper suscipit lobortis nisl ut aliquip ex ea commodo consequat. Duis autem vel eum iriure dolor in hendrerit in vulputate velit esse molestie consequat, vel illum dolore eu feugiat nulla facilisis at vero et accumsan et iusto odio dignissim qui blandit praesent luptatum zzril delenit augue duis dolore te feugait nulla facilisi. Lorem ipsum dolor sit amet, consectetuer adipiscing elit, sed diam nonummy nibh euismod tincidunt ut laoreet dolore magna aliquam erat volutpat. Ut wisi enim ad minim veniam, quis nostrud exerci tation ullamcorper suscipit lobortis nisl ut aliquip ex ea commodo consequat. Duis autem vel eum iriure dolor in hendrerit in vulputate velit esse molestie consequat, vel illum dolore eu feugiat nulla facilisis at vero et accumsan et iusto odio dignissim qui blandit praesent luptatum zzril delenit augue duis dolore te feugait nulla facilisi. Nam liber tempor cum soluta nobis eleifend option congue nihil imperdiet doming id quod mazim placerat facer possim assum.\\\\
\textbf{Schlagworte:} ABC, DEF
\end{abstract} 
 \cleardoublepage
\pagenumbering{roman}
\KOMAoptions{cleardoublepage=plain}
 \chapter*{Erklärung}

Hiermit erkl\"{a}re ich an Eides statt, dass ich die vorliegende Arbeit selbst\"{a}ndig und ohne 
fremde Hilfe verfasst, andere als die angegebenen Quellen und Hilfsmittel nicht benutzt 
und die aus anderen Quellen entnommenen Stellen als solche gekennzeichnet habe. 
\\\\\\
\line(1,0){200}\\
Furtwangen, \today, Schlauer Schreiberling
 \cleardoublepage

\tableofcontents % Inhaltsverzeichnis

\listoffigures % Abbildungsverzeichnis

\listoftables % Tabellenverzeichnis

\cleardoublepage

% Hauptinhalt; unterteilt in einzelne Kapitel; entsprechend auch einzelne Dateien.
% ********************************************************************
\pagenumbering{arabic}
\pagestyle{headings} % Kopfzeile mit lebendem Kolumnentitel
\parskip0.7ex % Abstand nach einem Absatz.
\chapter{Einleitung}
\section{Begründung der Themenstellung}

Lorem ipsum dolor sit amet, consectetuer adipiscing elit, sed diam nonummy nibh euismod tincidunt ut laoreet dolore magna aliquam erat volutpat. Ut wisi enim ad minim veniam, quis nostrud exerci tation ullamcorper suscipit lobortis nisl ut aliquip ex ea commodo consequat. Duis autem vel eum iriure dolor in hendrerit in vulputate velit esse molestie consequat, vel illum dolore eu feugiat nulla facilisis at vero et accumsan et iusto odio dignissim qui blandit praesent luptatum zzril delenit augue duis dolore te feugait nulla facilisi. Lorem ipsum dolor sit amet, consectetuer adipiscing elit, sed diam nonummy nibh euismod tincidunt ut laoreet dolore magna aliquam erat volutpat. Ut wisi enim ad minim veniam, quis nostrud exerci tation ullamcorper suscipit lobortis nisl ut aliquip ex ea commodo consequat. Duis autem vel eum iriure dolor in hendrerit in vulputate velit esse molestie consequat, vel illum dolore eu feugiat nulla facilisis at vero et accumsan et iusto odio dignissim qui blandit praesent luptatum zzril delenit augue duis dolore te feugait nulla facilisi. Nam liber tempor cum soluta nobis eleifend option congue nihil imperdiet doming id quod mazim placerat facer possim assum.

\section{Abgrenzung des Themas}

Lorem ipsum dolor sit amet, consectetuer adipiscing elit, sed diam nonummy nibh euismod tincidunt ut laoreet dolore magna aliquam erat volutpat. Ut wisi enim ad minim veniam, quis nostrud exerci tation ullamcorper suscipit lobortis nisl ut aliquip ex ea commodo consequat. Duis autem vel eum iriure dolor in hendrerit in vulputate velit esse molestie consequat, vel illum dolore eu feugiat nulla facilisis at vero et accumsan et iusto odio dignissim qui blandit praesent luptatum zzril delenit augue duis dolore te feugait nulla facilisi. Lorem ipsum dolor sit amet, consectetuer adipiscing elit, sed diam nonummy nibh euismod tincidunt ut laoreet dolore magna aliquam erat volutpat. Ut wisi enim ad minim veniam, quis nostrud exerci tation ullamcorper suscipit lobortis nisl ut aliquip ex ea commodo consequat. Duis autem vel eum iriure dolor in hendrerit in vulputate velit esse molestie consequat, vel illum dolore eu feugiat nulla facilisis at vero et accumsan et iusto odio dignissim qui blandit praesent luptatum zzril delenit augue duis dolore te feugait nulla facilisi. Nam liber tempor cum soluta nobis eleifend option congue nihil imperdiet doming id quod mazim placerat facer possim assum.
 
\section{Aufbau und Struktur der Arbeit}

Lorem ipsum dolor sit amet, consectetuer adipiscing elit, sed diam nonummy nibh euismod tincidunt ut laoreet dolore magna aliquam erat volutpat. Ut wisi enim ad minim veniam, quis nostrud exerci tation ullamcorper suscipit lobortis nisl ut aliquip ex ea commodo consequat. Duis autem vel eum iriure dolor in hendrerit in vulputate velit esse molestie consequat, vel illum dolore eu feugiat nulla facilisis at vero et accumsan et iusto odio dignissim qui blandit praesent luptatum zzril delenit augue duis dolore te feugait nulla facilisi. Lorem ipsum dolor sit amet, consectetuer adipiscing elit, sed diam nonummy nibh euismod tincidunt ut laoreet dolore magna aliquam erat volutpat. Ut wisi enim ad minim veniam, quis nostrud exerci tation ullamcorper suscipit lobortis nisl ut aliquip ex ea commodo consequat. Duis autem vel eum iriure dolor in hendrerit in vulputate velit esse molestie consequat, vel illum dolore eu feugiat nulla facilisis at vero et accumsan et iusto odio dignissim qui blandit praesent luptatum zzril delenit augue duis dolore te feugait nulla facilisi. Nam liber tempor cum soluta nobis eleifend option congue nihil imperdiet doming id quod mazim placerat facer possim assum.
\chapter{Überschrift erster Ordnung}


Lorem ipsum dolor sit amet, consectetuer adipiscing elit, sed diam nonummy nibh euismod tincidunt ut laoreet dolore magna aliquam erat volutpat. Ut wisi enim ad minim veniam, quis nostrud exerci tation ullamcorper suscipit lobortis nisl ut aliquip ex ea commodo consequat. Duis autem vel eum iriure dolor in hendrerit in vulputate velit esse molestie consequat, vel illum dolore eu feugiat nulla facilisis at vero et accumsan et iusto odio dignissim qui blandit praesent luptatum zzril delenit augue duis dolore te feugait nulla facilisi. Lorem ipsum dolor sit amet, consectetuer adipiscing elit, sed diam nonummy nibh euismod tincidunt ut laoreet dolore magna aliquam erat volutpat. Ut wisi enim ad minim veniam, quis nostrud exerci tation ullamcorper suscipit lobortis nisl ut aliquip ex ea commodo consequat. Duis autem vel eum iriure dolor in hendrerit in vulputate velit esse molestie consequat, vel illum dolore eu feugiat nulla facilisis at vero et accumsan et iusto odio dignissim qui blandit praesent luptatum zzril delenit augue duis dolore te feugait nulla facilisi. Nam liber tempor cum soluta nobis eleifend option congue nihil imperdiet doming id quod mazim placerat facer possim assum.

\section{Überschrift zweiter Ordnung}\label{bsp}

Lorem ipsum dolor sit amet, consectetuer adipiscing elit, sed diam nonummy nibh euismod tincidunt ut laoreet dolore magna aliquam erat volutpat. Ut wisi enim ad minim veniam, quis nostrud exerci tation ullamcorper suscipit lobortis nisl ut aliquip ex ea commodo consequat. Duis autem vel eum iriure dolor in hendrerit in vulputate velit esse molestie consequat, vel illum dolore eu feugiat nulla facilisis at vero et accumsan et iusto odio dignissim qui blandit praesent luptatum zzril delenit augue duis dolore te feugait nulla facilisi. Lorem ipsum dolor sit amet, consectetuer adipiscing elit, sed diam nonummy nibh euismod tincidunt ut laoreet dolore magna aliquam erat volutpat. Ut wisi enim ad minim veniam, quis nostrud exerci tation ullamcorper suscipit lobortis nisl ut aliquip ex ea commodo consequat. Duis autem vel eum iriure dolor in hendrerit in vulputate velit esse molestie consequat, vel illum dolore eu feugiat nulla facilisis at vero et accumsan et iusto odio dignissim qui blandit praesent luptatum zzril delenit augue duis dolore te feugait nulla facilisi. Nam liber tempor cum soluta nobis eleifend option congue nihil imperdiet doming id quod mazim placerat facer possim assum.

\subsection{Überschrift dritter Ordnung}

Hier steht Text mit einem Verweis auf Kapitel \ref{bsp}. 

\section{Und weiter}

Und hier steht noch mehr Text. Und hier ein Literaturverweis: \cite{bortz-2006}. Toll, oder?


% Appendix für Grafiken und Evaluationsfragen etc.
% ********************************************************************
 \appendix
  \chapter{Appendix A}\label{apx:A}
  \chapter{Appendix B}\label{apx:B}
  \chapter{Appendix C}\label{apx:C}



% Literaturverzeichnis
% ********************************************************************

\bibliographystyle{dcu_url}
\clearpage
\addcontentsline{toc}{chapter}{Literaturverzeichnis}
\bibliography{Literatursammlung} %BibTeX-Literaturdatei

% \addcontentsline{toc}{chapter}{Index}
% \printindex
\end{document}
